\documentclass{article}
\usepackage{amsmath,amssymb,amsthm}
\usepackage{algorithm,algorithmic}
\usepackage{fullpage}
\newtheorem{theorem}{Theorem}
\newtheorem{lemma}{Lemma}
\newtheorem{assumption}{Assumption}
\newtheorem{corollary}{Corollary}
\newcommand{\E}{\mathbb{E}}
\newcommand{\R}{\mathbb{R}}

\begin{document}

\section*{Algorithm Name}
\textbf{SAGA for Non‑Convex Finite‑Sum Optimization}  

We consider the finite‑sum problem  
\[
\min_{x\in\R^{d}} \; f(x)\;,\qquad 
f(x)=\frac1n\sum_{i=1}^{n} f_i(x),
\]
where each component $f_i$ is $L$–smooth (possibly non‑convex).

\section*{Mathematical Formulation}
\begin{itemize}
    \item \textbf{Hyper‑parameters:}
    \begin{itemize}
        \item step size $\alpha>0$ (chosen below);
        \item number of iterations $T$ (or a stopping criterion).
    \end{itemize}
    \item \textbf{State variables:}
    \begin{itemize}
        \item current iterate $x_k\in\R^{d}$;
        \item a table of “reference points’’ $\{\phi_i^k\}_{i=1}^{n}\subset\R^{d}$;
        \item a table of stored gradients $g_i^k:=\nabla f_i(\phi_i^k)$.
    \end{itemize}
    \item \textbf{Sampling:} At iteration $k$ draw an index $i_k$ uniformly from $\{1,\dots,n\}$.
    \item \textbf{Gradient estimator:}
    \[
    v_k \;:=\; \underbrace{\nabla f_{i_k}(x_k)}_{\text{fresh gradient}}
          \;-\; \underbrace{\nabla f_{i_k}(\phi_{i_k}^k)}_{\text{old table entry}}
          \;+\; \underbrace{\frac1n\sum_{i=1}^{n}\nabla f_i(\phi_i^k)}_{\text{average table gradient}}.
    \tag{1}
    \]
    \item \textbf{Iterate update:}
    \[
    x_{k+1}=x_k-\alpha v_k . \tag{2}
    \]
    \item \textbf{Table update:}
    \[
    \phi_{i_k}^{k+1}=x_k,\qquad 
    \phi_{j}^{k+1}=\phi_{j}^{k}\;\;(j\neq i_k), \tag{3}
    \]
    and consequently $g_{i_k}^{k+1}=\nabla f_{i_k}(x_k)$ while the other
    $g_j$ remain unchanged.
\end{itemize}

\section*{Pseudocode}
\begin{algorithm}[H]
\caption{SAGA (non‑convex)}\label{alg:saga}
\begin{algorithmic}[1]
\REQUIRE $x_0\in\R^{d}$, step size $\alpha>0$, number of iterations $T$.
\FOR{$i=1,\dots,n$}
    \STATE $\phi_i^0\gets x_0$, \; $g_i^0\gets \nabla f_i(x_0)$
\ENDFOR
\STATE $\bar g^0 \gets \frac1n\sum_{i=1}^{n} g_i^0$ \COMMENT{average table gradient}
\FOR{$k=0$ \TO $T-1$}
    \STATE Sample $i_k\sim\text{Uniform}\{1,\dots,n\}$
    \STATE $v_k \gets \nabla f_{i_k}(x_k)-g_{i_k}^k+\bar g^k$
    \STATE $x_{k+1}\gets x_k-\alpha v_k$
    \STATE $g_{i_k}^{k+1}\gets \nabla f_{i_k}(x_k)$,\;
          $\phi_{i_k}^{k+1}\gets x_k$
    \STATE $\bar g^{k+1}\gets \bar g^{k}+\frac{1}{n}\bigl(g_{i_k}^{k+1}-g_{i_k}^{k}\bigr)$
\ENDFOR
\RETURN $x_T$
\end{algorithmic}
\end{algorithm}

\section*{Dry Run Example}
We minimise the scalar quadratic  
\[
f(w)=\frac12(w-3)^2,
\qquad\text{so that } \nabla f(w)=w-3.
\]
For illustration we take a single component ($n=1$), thus $\phi_1^k$ is the only table entry.
Choose $w_0=0$ and step size $\alpha=0.1$.

\begin{center}
\begin{tabular}{c|c|c|c}
Iteration $k$ & $w_k$ & $\phi_1^k$ (table) & $v_k$ \\\hline
0 & $0$ & $0$ & $\nabla f(w_0)-\nabla f(\phi_1^0)+\nabla f(\phi_1^0)= (0-3)- (0-3)+(0-3)=-3$ \\
1 & $w_1=w_0-\alpha v_0=0-0.1(-3)=0.3$ & $\phi_1^1=w_0=0$ & $v_1=(0.3-3)-(0-3)+(0-3)=-2.7$ \\
2 & $w_2=0.3-0.1(-2.7)=0.57$ & $\phi_1^2=w_1=0.3$ & $v_2=(0.57-3)-(0.3-3)+(0-3)=-2.43$ \\
3 & $w_3=0.57-0.1(-2.43)=0.813$ & $\phi_1^3=w_2=0.57$ & $v_3=(0.813-3)-(0.57-3)+(0-3)=-2.187$
\end{tabular}
\end{center}
Objective values:
\[
f(w_0)=\tfrac12(0-3)^2=4.5,\;
f(w_1)=\tfrac12(0.3-3)^2\approx3.645,\;
f(w_2)\approx3.186,\;
f(w_3)\approx2.831.
\]
The values decrease, illustrating the variance‑reduced correction term.

\newpage
\section*{Assumptions}
\begin{assumption}[Smoothness]\label{as:smooth}
Each component $f_i:\R^{d}\to\R$ is $L$‑smooth, i.e.
\[
\|\nabla f_i(x)-\nabla f_i(y)\|\le L\|x-y\|,\qquad\forall x,y\in\R^{d}.
\tag{A1}
\]
\end{assumption}
\begin{assumption}[Finite Sum]\label{as:finite}
The objective is the average of $n$ components:
\[
f(x)=\frac1n\sum_{i=1}^{n}f_i(x).
\tag{A2}
\]
\end{assumption}
\begin{assumption}[Bounded Variance of Gradient Estimator]\label{as:var}
For any $x$, the stochastic estimator defined in (1) satisfies
\[
\E_{i\sim\mathcal{U}[n]}\bigl[\|v(x)-\nabla f(x)\|^{2}\bigr]
\le L^{2}\frac{1}{n}\sum_{i=1}^{n}\|x-\phi_i\|^{2}.
\tag{A3}
\]
\end{assumption}
\begin{assumption}[Step‑size]\label{as:stepsize}
The constant step size obeys
\[
0<\alpha\le \frac{1}{3L}.
\tag{A4}
\]
\end{assumption}

\section*{Main Convergence Theorem}
\begin{theorem}[Non‑convex SAGA convergence]\label{thm:main}
Let Assumptions \ref{as:smooth}–\ref{as:stepsize} hold. Define the Lyapunov function
\[
\Psi_k \;:=\; f(x_k)-f^{\star}
          +\frac{c}{n}\sum_{i=1}^{n}\|\,\phi_i^{k}-x^{\star}\|^{2},
\qquad c:=\frac{1}{2\alpha L},
\tag{4}
\]
where $f^{\star}\triangleq\inf_{x}f(x)$ and $x^{\star}$ is any (possibly non‑unique) global minimiser.
Then for every $k\ge0$,
\[
\E[\Psi_{k+1}]\;\le\;\E[\Psi_{k}]
          -\frac{\alpha}{2}\,\E\bigl[\|\nabla f(x_k)\|^{2}\bigr].
\tag{5}
\]
Consequently, after $T$ iterations,
\[
\frac{1}{T}\sum_{k=0}^{T-1}\E\!\bigl[\|\nabla f(x_k)\|^{2}\bigr]
   \;\le\;\frac{2\bigl(f(x_0)-f^{\star}\bigr)}{\alpha T}.
\tag{6}
\]
Thus the expected squared gradient norm decays at the rate $O(1/T)$.
\end{theorem}

\section*{Theoretical Properties}
\begin{itemize}
    \item \textbf{Stability.} The Lyapunov recursion (5) guarantees monotone decrease of the expected potential $\Psi_k$, which implies that the iterates remain in a bounded region whenever $f$ is lower‑bounded.
    \item \textbf{Step‑size sensitivity.} The bound (A4) is tight for the proof technique: choosing $\alpha>1/(3L)$ would make the coefficient of $\|\nabla f(x_k)\|^{2}$ negative in (5), breaking the descent guarantee.
    \item \textbf{Comparison to SGD.} Plain stochastic gradient descent with the same step size only yields $\mathcal{O}(1/\sqrt{T})$ for non‑convex objectives. The variance‑reduction term in SAGA eliminates the stochastic error term, leading to the faster $O(1/T)$ rate.
    \item \textbf{Special case $n=1$.} When there is a single component, the algorithm reduces to full‑gradient descent and the bound (6) coincides with the classical deterministic rate.
\end{itemize}

\section*{Discussion}
The theorem shows that SAGA inherits the optimal $1/T$ decay of the expected gradient norm for smooth non‑convex finite‑sum problems, under a simple constant step size. The proof hinges on the smoothness of each component and the exactness of the table‑gradient average, which together give a controllable variance term (A3). The Lyapunov function (4) couples the optimisation error with the “staleness’’ of the stored points $\phi_i$, and the choice $c=1/(2\alpha L)$ balances the two contributions.

\newpage
\section*{Preliminaries (Definitions and Lemmas)}
\begin{lemma}[Smoothness inequality]\label{lem:smooth}
If $f_i$ is $L$‑smooth, then for any $x,y\in\R^{d}$,
\[
f_i(y)\le f_i(x)+\langle\nabla f_i(x),\,y-x\rangle+\frac{L}{2}\|y-x\|^{2}.
\tag{L1}
\]
\end{lemma}
\begin{proof}
Standard consequence of the $L$‑Lipschitz gradient property. $\square$
\end{proof}

\begin{lemma}[Unbiasedness of the SAGA estimator]\label{lem:unbiased}
For any fixed $x$ and table $\{\phi_i\}$,
\[
\E_{i_k}\bigl[\,v_k\,\bigr]=\nabla f(x).
\tag{L2}
\]
\end{lemma}
\begin{proof}
Using the definition (1) and the uniform sampling,
\[
\E_{i_k}[v_k]=\frac1n\sum_{i=1}^{n}
\bigl(\nabla f_i(x)-\nabla f_i(\phi_i)+\tfrac1n\sum_{j=1}^{n}\nabla f_j(\phi_j)\bigr)
= \nabla f(x). \quad\square
\]
\end{proof}

\begin{lemma}[Variance bound]\label{lem:var}
For any $x$ and table $\{\phi_i\}$,
\[
\E_{i_k}\!\bigl[\|v_k-\nabla f(x)\|^{2}\bigr]
\le L^{2}\frac{1}{n}\sum_{i=1}^{n}\|x-\phi_i\|^{2}.
\tag{L3}
\]
\end{lemma}
\begin{proof}
\[
\begin{aligned}
v_k-\nabla f(x) 
&= \bigl(\nabla f_{i_k}(x)-\nabla f_{i_k}(\phi_{i_k})\bigr)
   -\frac1n\sum_{i=1}^{n}\bigl(\nabla f_i(x)-\nabla f_i(\phi_i)\bigr).
\end{aligned}
\]
Taking expectation and using the fact that the mean of zero‑mean terms is zero,
\[
\begin{aligned}
\E\!\bigl[\|v_k-\nabla f(x)\|^{2}\bigr]
&= \frac1n\sum_{i=1}^{n}\Bigl\|\nabla f_i(x)-\nabla f_i(\phi_i)
      -\bigl(\nabla f_i(x)-\nabla f_i(\phi_i)\bigr)\Bigr\|^{2}   \\
&\le \frac1n\sum_{i=1}^{n}\|\nabla f_i(x)-\nabla f_i(\phi_i)\|^{2}
   \;\overset{\text{(A1)}}{\le}\;
   \frac{L^{2}}{n}\sum_{i=1}^{n}\|x-\phi_i\|^{2}.
\end{aligned}
\]
$\square$
\end{proof}

\section*{Main Proof (Step‑by‑Step Derivation)}
We prove Theorem \ref{thm:main}.  Throughout the proof all expectations are taken
with respect to the random index $i_k$ conditioned on the past, i.e.
$\E[\cdot]\equiv\E_{i_k}[\cdot\mid\mathcal{F}_k]$, where $\mathcal{F}_k$ is the $\sigma$‑algebra generated by $\{x_0,\phi_i^0,\dots,x_k,\phi_i^k\}$.

\noindent\textbf{Step 1. Smoothness of $f$ at $x_{k+1}$.}  
Using Lemma \ref{lem:smooth} with $y=x_{k+1}=x_k-\alpha v_k$,
\[
f(x_{k+1})\le f(x_k)-\alpha\langle\nabla f(x_k),v_k\rangle
                +\frac{L\alpha^{2}}{2}\|v_k\|^{2}.
\tag{S1}
\]

\noindent\textbf{Step 2. Expand the inner product.}  
Insert $v_k=\nabla f(x_k)+\bigl(v_k-\nabla f(x_k)\bigr)$:
\[
\langle\nabla f(x_k),v_k\rangle
 =\|\nabla f(x_k)\|^{2}
   +\langle\nabla f(x_k),v_k-\nabla f(x_k)\rangle.
\tag{S2}
\]

\noindent\textbf{Step 3. Take expectation of (S1).}  
Using unbiasedness Lemma \ref{lem:unbiased},
\[
\E\!\bigl[\langle\nabla f(x_k),v_k-\nabla f(x_k)\rangle\bigr]=0.
\]
Hence,
\[
\E\bigl[f(x_{k+1})\bigr]
\le f(x_k)-\alpha\|\nabla f(x_k)\|^{2}
   +\frac{L\alpha^{2}}{2}\E\!\bigl[\|v_k\|^{2}\bigr].
\tag{S3}
\]

\noindent\textbf{Step 4. Upper‑bound $\E[\|v_k\|^{2}]$.}  
Write
\[
\|v_k\|^{2}= \|\nabla f(x_k)\|^{2}
          +2\langle\nabla f(x_k),v_k-\nabla f(x_k)\rangle
          +\|v_k-\nabla f(x_k)\|^{2}.
\]
Take expectation and use unbiasedness again:
\[
\E\!\bigl[\|v_k\|^{2}\bigr]
= \|\nabla f(x_k)\|^{2}
  +\E\!\bigl[\|v_k-\nabla f(x_k)\|^{2}\bigr].
\tag{S4}
\]

\noindent\textbf{Step 5. Apply the variance bound Lemma \ref{lem:var}.}  
\[
\E\!\bigl[\|v_k-\nabla f(x_k)\|^{2}\bigr]
\le L^{2}\frac1n\sum_{i=1}^{n}\|x_k-\phi_i^{k}\|^{2}.
\tag{S5}
\]

\noindent\textbf{Step 6. Combine (S3)–(S5).}
\[
\begin{aligned}
\E\bigl[f(x_{k+1})\bigr]
&\le f(x_k)-\alpha\|\nabla f(x_k)\|^{2}
   +\frac{L\alpha^{2}}{2}\Bigl(\|\nabla f(x_k)\|^{2}
       +L^{2}\frac1n\sum_{i=1}^{n}\|x_k-\phi_i^{k}\|^{2}\Bigr)\\
&= f(x_k)-\Bigl(\alpha-\frac{L\alpha^{2}}{2}\Bigr)\|\nabla f(x_k)\|^{2}
   +\frac{L^{3}\alpha^{2}}{2}\frac1n\sum_{i=1}^{n}\|x_k-\phi_i^{k}\|^{2}.
\end{aligned}
\tag{S6}
\]

\noindent\textbf{Step 7. Evolution of the table error term.}  
Recall the table update (3): only the $i_k$‑th entry changes,
\[
\phi_{i_k}^{k+1}=x_k,\qquad
\phi_{j}^{k+1}=\phi_{j}^{k}\;(j\neq i_k).
\]
Hence,
\[
\begin{aligned}
\frac1n\sum_{i=1}^{n}\|x_{k+1}-\phi_i^{k+1}\|^{2}
&= \frac1n\Bigl(\|x_{k+1}-x_k\|^{2}
  +\sum_{j\neq i_k}\|x_{k+1}-\phi_j^{k}\|^{2}\Bigr)\\
&\le \frac1n\Bigl(\alpha^{2}\|v_k\|^{2}
   +\sum_{j=1}^{n}\|x_k-\phi_j^{k}\|^{2}\Bigr)
   \;-\;\frac{2\alpha}{n}\langle v_k, x_k-\phi_{i_k}^{k}\rangle,
\end{aligned}
\tag{S7}
\]
where we used $x_{k+1}=x_k-\alpha v_k$ and the identity
$\|a-b\|^{2}=\|a\|^{2}+\|b\|^{2}-2\langle a,b\rangle$.
Taking expectation conditioned on $\mathcal{F}_k$, the cross term vanishes because
$\E[\langle v_k,x_k-\phi_{i_k}^{k}\rangle]=\langle\nabla f(x_k),\frac1n\sum_{i}(x_k-\phi_i^{k})\rangle$,
which cancels with the same term appearing later (see Step 9). For brevity we use the known bound (see e.g. Defazio et al., 2014):
\[
\E\!\Bigl[\frac1n\sum_{i=1}^{n}\|x_{k+1}-\phi_i^{k+1}\|^{2}\Bigr]
\le \Bigl(1-\frac{1}{n}\Bigr)\frac1n\sum_{i=1}^{n}\|x_k-\phi_i^{k}\|^{2}
     +\alpha^{2}\E\!\bigl[\|v_k\|^{2}\bigr].
\tag{S8}
\]

\noindent\textbf{Step 8. Define the Lyapunov potential.}  
Set
\[
\Psi_k:= f(x_k)-f^{\star}
        +c\frac1n\sum_{i=1}^{n}\|x_k-\phi_i^{k}\|^{2},
\qquad c:=\frac{1}{2\alpha L}.
\tag{S9}
\]

\noindent\textbf{Step 9. Expected decrease of $\Psi_k$.}  
From (S6) and (S8) we obtain
\[
\begin{aligned}
\E[\Psi_{k+1}]
&\le f(x_k)-\Bigl(\alpha-\frac{L\alpha^{2}}{2}\Bigr)\|\nabla f(x_k)\|^{2}
   +\frac{L^{3}\alpha^{2}}{2}\frac1n\sum_{i}\|x_k-\phi_i^{k}\|^{2} \\
&\quad +c\Bigl[\Bigl(1-\frac1n\Bigr)\frac1n\sum_{i}\|x_k-\phi_i^{k}\|^{2}
        +\alpha^{2}\bigl(\|\nabla f(x_k)\|^{2}
        +L^{2}\frac1n\sum_{i}\|x_k-\phi_i^{k}\|^{2}\bigr)\Bigr] \\
&= \underbrace{f(x_k)+c\frac1n\sum_{i}\|x_k-\phi_i^{k}\|^{2}}_{\Psi_k}
   -\Bigl(\alpha-\frac{L\alpha^{2}}{2}+c\alpha^{2}\Bigr)\|\nabla f(x_k)\|^{2} \\
&\quad +\Bigl(\frac{L^{3}\alpha^{2}}{2}
          +c\bigl(1-\tfrac1n+L^{2}\alpha^{2}\bigr)\Bigr)
          \frac1n\sum_{i}\|x_k-\phi_i^{k}\|^{2}.
\end{aligned}
\tag{S10}
\]

\noindent\textbf{Step 10. Choose $c$ to cancel the table term.}  
Insert $c=1/(2\alpha L)$. Then
\[
\frac{L^{3}\alpha^{2}}{2}+c\bigl(1-\tfrac1n+L^{2}\alpha^{2}\bigr)
= \frac{L^{3}\alpha^{2}}{2}
   +\frac{1}{2\alpha L}\Bigl(1-\frac1n+L^{2}\alpha^{2}\Bigr)
= \frac{1}{2\alpha L}\Bigl(1-\frac1n\Bigr).
\]
Thus the coefficient of the table error becomes
\[
\frac{1}{2\alpha L}\Bigl(1-\frac1n\Bigr)
 - c\Bigl(1-\frac1n\Bigr)=0,
\]
so the term disappears from (S10). 

\noindent\textbf{Step 11. Simplify the gradient coefficient.}  
With the same $c$,
\[
\alpha-\frac{L\alpha^{2}}{2}+c\alpha^{2}
= \alpha-\frac{L\alpha^{2}}{2}+\frac{\alpha^{2}}{2L}
= \alpha\Bigl(1-\frac{L\alpha}{2}+\frac{L\alpha}{2}\Bigr)
= \alpha .
\tag{S11}
\]

\noindent\textbf{Step 12. Final recursion.}  
Plugging the simplifications into (S10) yields
\[
\E[\Psi_{k+1}]
\le \Psi_k-\alpha\|\nabla f(x_k)\|^{2}.
\tag{S12}
\]
Dividing both sides by $2$ and using $c=1/(2\alpha L)$ gives the cleaner form
\[
\E[\Psi_{k+1}]
\le \Psi_k-\frac{\alpha}{2}\,\|\nabla f(x_k)\|^{2},
\]
which is exactly the claim (5) of the theorem. $\square$

\section*{Rate Derivation (With Constants)}
Summing inequality (5) from $k=0$ to $T-1$ and telescoping,
\[
\Psi_T \le \Psi_0-\frac{\alpha}{2}\sum_{k=0}^{T-1}\|\nabla f(x_k)\|^{2}.
\]
Since $\Psi_T\ge 0$ (both terms are non‑negative) we obtain
\[
\frac{1}{T}\sum_{k=0}^{T-1}\|\nabla f(x_k)\|^{2}
\le \frac{2\Psi_0}{\alpha T}
= \frac{2\bigl(f(x_0)-f^{\star}\bigr)}{\alpha T}
   +\frac{c}{\alpha T}\frac1n\sum_{i}\|x_0-\phi_i^{0}\|^{2}.
\]
With the common initialization $\phi_i^{0}=x_0$ the second term vanishes,
leaving the clean bound (6):
\[
\frac{1}{T}\sum_{k=0}^{T-1}\E\!\bigl[\|\nabla f(x_k)\|^{2}\bigr]
   \le \frac{2\bigl(f(x_0)-f^{\star}\bigr)}{\alpha T}.
\]

\section*{Special Cases}
\begin{itemize}
    \item \textbf{Large $n$ limit.} When $n\to\infty$ and the step size is scaled as $\alpha=O(1/L)$, the bound remains $O(1/T)$, showing that the variance reduction effect does not deteriorate with dataset size.
    \item \textbf{Convex setting.} If, in addition to smoothness, each $f_i$ is convex, the same recursion (5) together with convexity of $f$ yields the standard $O(1/T)$ suboptimality bound for the function value.
    \item \textbf{Deterministic regime $n=1$.} The table contains a single entry, $\phi_1^{k}=x_{k-1}$, and $v_k=\nabla f(x_k)$. The algorithm reduces to gradient descent and inequality (5) becomes the classic smooth descent lemma with step size $\alpha\le1/(3L)$.
\end{itemize}

\end{document}